% !TeX program = lualatex
% !TeX encoding = utf8
% !TeX spellcheck = uk_UA
% !TeX root = ../main.tex

%=========================================================
\chapter{Елементи фізики плазми}\label{\currfilebase}\hypertarget{PlasmaPhysics}{}
\graphicspath{{\currfilebase/Pictures}}
%=========================================================

У цьому розділі використовується гаусівська абсолютна система одиниць (СГС). Вибір цієї системи у фізиці плазми є загальноприйнятим стандартом, оскільки вона дозволяє уникнути штучних коефіцієнтів $\varepsilon_0$ та $\mu_0$, явно підкреслюючи симетрію електричного та магнітного полів і фундаментальну роль швидкості світла $c$.

У системах СГС:
\begin{itemize}
	\item Сила Лоренца має вигляд: $\vect{F} = q\left(\vect{E} + \frac{1}{c}\vect{v}\times\vect{B}\right)$;
	\item Рівняння Пуассона записується як: $\nabla^2\varphi = -4\pi\rho$;
	\item Густина енергії магнітного поля становить: $w_B = \frac{B^2}{8\pi}$.
\end{itemize}

Температура $T$ всюди вказує термодинамічну температуру, тому середня теплова енергія на один ступінь вільності дорівнює $\frac{1}{2}k_{\mathrm B}T$, де $k_{\mathrm B} = 1.38\times10^{-16}~\text{ерг/К}$~--- стала Больцмана. У фізиці плазми температуру традиційно вимірюють в електронвольтах ($\mathrm{eV}$), що відповідає енергії $k_{\mathrm B}T$. Зв'язок між одиницями має вигляд:
\begin{equation}
	1~\mathrm{eV} \leftrightarrow \frac{1~\mathrm{eV}}{k_{\mathrm B}} \simeq 1.160\times10^4~\mathrm{K}.
\end{equation}
Наприклад, термоядерна плазма з температурою $T = 10~\mathrm{keV}$ відповідає приблизно $1.16\times10^8~\mathrm{K}$.

%% ========================================================
\section{Вступ. Що таке плазма}\label{sec:plasma-basics}\hypertarget{plasma-basics}{}
%% ========================================================

%% --------------------------------------------------------
\subsection{Методичний вступ: від нейтрального газу до плазми}
%% --------------------------------------------------------

\textbf{Чому виникає потреба в окремому розділі фізики?}
У звичайному нейтральному газі взаємодія між молекулами має короткодіючий (ван-дер-ваальсівський) характер. Частинки рухаються прямолінійно й взаємодіють лише під час короткочасних парних зіткнень.

Однак, якщо нагріти газ або піддати його сильному електромагнітному випромінюванню, відбувається іонізація: атоми розпадаються на вільні електрони та позитивні іони. Внаслідок цього виникає принципово новий стан речовини~--- \emph{четвертий стан речовини}, або \emph{плазма}.

Головна відмінність плазми від нейтрального газу полягає в наявності \textbf{кулонівських сил}, які падають із відстанню як $1/r^2$. Кулонівська взаємодія є далекодіючою: кожна заряджена частинка одночасно взаємодіє з величезною кількістю навколишніх частинок. Це породжує так звану \emph{колективну поведінку}~--- здатність плазми реагувати на зовнішні та внутрішні збурення не як набір ізольованих частинок, а як єдине узгоджене суцільне середовище.

%% --------------------------------------------------------
\subsection{Визначення плазми та ступінь іонізації}
%% --------------------------------------------------------

\emphx{Плазма} (від грец. \emphx{plasma}~--- виліплений, оформлений)~--- це частково або повністю іонізований газ, в якому густини позитивних і негативних зарядів практично однакові (квазінейтральність), а кулонівська взаємодія визначає його динаміку та колективні властивості. Термін запропонований у 1929 році Ірвінгом Ленгмюром.

Для кількісної характеристики стану іонізації вводиться \emphx{ступінь іонізації} $\alpha$:
\begin{equation}
	\alpha = \frac{n_i}{n_i + n_n},
\end{equation}
де $n_i$~--- об'ємна концентрація іонів (число частинок в $1~\mathrm{см}^3$), а $n_n$~--- концентрація нейтральних атомів або молекул.
\begin{itemize}
	\item Якщо $\alpha \ll 1$ (наприклад, $\alpha \sim 10^{-6}$--$10^{-4}$ у слабоіонізованому розряді), плазму називають \emph{слабоіонізованою}.
	\item Якщо $\alpha \to 1$ ($n_n \ll n_i$), плазма є \emph{повністю іонізованою} (зоряна речовина, термоядерний реактор).
\end{itemize}

За температурним масштабом плазму поділяють на:
\begin{itemize}
	\item \emphx{Низькотемпературну} ($T_e < 10^5~\mathrm{K} \approx 10~\mathrm{eV}$): газові розряди, іоносфера, плазмові технології;
	\item \emphx{Високотемпературну} ($T_e > 10^6$--$10^8~\mathrm{K} \approx 0.1$--$10~\mathrm{keV}$): ядра зірок, експериментальні установки керованого термоядерного синтезу (КТС).
\end{itemize}

%% --------------------------------------------------------
\subsection{Двокомпонентність та багатотемпературність плазми}
%% --------------------------------------------------------

\textbf{Фізична проблема:} Чому плазму не завжди можна описати однією термодинамічною температурою?

Плазма складається з частинок з колосальною різницею мас. Маса найлегшого іона (протона) $m_p$ перевищує масу електрона $m_e$ у 1836 разів:
\begin{equation}
	\frac{m_e}{m_i} \ll 1 \quad \left(\text{для дейтерію } \frac{m_e}{m_d} \approx \frac{1}{3670}\right).
\end{equation}

Під час пружного зіткнення легкої частинки ($m_e$) з важкою ($m_i$) частка переданої кінетичної енергії за один удар дорівнює за порядком величини:
\begin{equation}
	\Delta E \sim \frac{m_e}{m_i} E_e \ll E_e.
\end{equation}
Тому енергообмін між електронною та іонною компонентами відбувається дуже повільно. У газорозрядній плазмі зовнішнє електричне поле передає енергію насамперед легким електронам. Вони швидко досягають стану теплової рівноваги між собою із температурою $T_e$, але не встигають передати цю енергію іонам.

Внаслідок цього виникає \emphx{неізотермічна (двотемпературна) плазма}, в якій:
\begin{equation}
	T_e \gg T_i.
\end{equation}

Часи встановлення внутрішньої рівноваги (часи релаксації) підпорядковуються співвідношенню:
\begin{equation}\label{eq:relaxation-times-ratio}
	\tau_e : \tau_i : \tau_{ie} \sim 1 : \sqrt{\frac{m_i}{m_e}} : \frac{m_i}{m_e},
\end{equation}
де $\tau_e$~--- час встановлення максвеллiвського розподілу серед електронів, $\tau_i$~--- серед іонів, а $\tau_{ie}$~--- час вирівнювання температур між електронами та іонами.

%% ========================================================
\section{Просторові масштаби та екранування в плазмі}\label{sec:spatial-scales}
%% ========================================================

%% --------------------------------------------------------
\subsection{Квазінейтральність плазми}
%% --------------------------------------------------------

\textbf{Педагогічна мотивація:} Що станеться, якщо в суцільному об'ємі плазми створити нескомпенсований заряд?

Нехай у плазмі з концентраціями $n_e$ електронів (заряд $-e$) та $n_i$ іонів (заряд $+Ze$) виникла просторова локальна різниця густин. За законом Кулона виникне кулонівське поле:
\begin{equation}
	\nabla \cdot \vect{E} = 4\pi \rho = 4\pi e (Z n_i - n_e).
\end{equation}
Оскільки питомий заряд електрона $e/m_e$ величезний, це поле миттєво створить потужні сили, які змістять електрони так, щоб знівелювати розділення зарядів.

Тому на макроскопічних масштабах плазма є \emphx{квазінейтральною}:
\begin{equation}\label{eq:quasineutrality}
	\rho = \sum_s q_s n_s = e(Z n_i - n_e) \approx 0 \quad \Rightarrow \quad n_e \approx Z n_i.
\end{equation}
Квазінейтральність~--- це не тотожний нуль у кожній точці, а мікроскопічна динамічна рівновага. Відхилення від неї можливі лише на дуже малих просторових масштабах.

%% --------------------------------------------------------
\subsection{Дебаївське екранування та електронний радіус Дебая}
%% --------------------------------------------------------

\textbf{Фізична проблема:} На якій саме відстані плазма здатна замаскувати (заекранувати) внесений у неї зовнішній точковий заряд $q$?

Розглянемо пробний заряд $+q$, вміщений у точку $r = 0$ суцільної плазми. Електрони з температурою $T_e$ притягуються до заряду, а іони~--- відштовхуються. Розподіл концентрації електронів у потенціалі $\varphi(r)$ описується статистикою Больцмана:
\begin{equation}
	n_e(r) = n_0 \exp\left(\frac{e\varphi(r)}{k_{\mathrm B}T_e}\right),
\end{equation}
де $n_0$~--- незбурена концентрація на нескінченності (де $\varphi \to 0$), $e > 0$~--- елементарний заряд.

У наближенні слабкої взаємодії ($e\varphi \ll k_{\mathrm B}T_e$) розкладемо експоненту у ряд Тейлора:
\begin{equation}
	n_e(r) \approx n_0 \left(1 + \frac{e\varphi}{k_{\mathrm B}T_e}\right).
\end{equation}

Вважаючи іони нерухомими усередненими носіями фону ($n_i = n_0/Z$), обчислимо густину заряду плазми $\rho_{\mathrm{pl}}$:
\begin{equation}
	\rho_{\mathrm{pl}} = Z e n_i - e n_e \approx e n_0 - e n_0\left(1 + \frac{e\varphi}{k_{\mathrm B}T_e}\right) = -\frac{n_0 e^2}{k_{\mathrm B}T_e}\varphi.
\end{equation}

Підставимо $\rho_{\mathrm{pl}}$ у рівняння Пуассона для радіально-симетричного поля при $r > 0$:
\begin{equation}
	\nabla^2\varphi = -4\pi\rho_{\mathrm{pl}} \quad \Rightarrow \quad \frac{1}{r^2}\frac{d}{dr}\left(r^2\frac{d\varphi}{dr}\right) = \frac{4\pi n_0 e^2}{k_{\mathrm B}T_e}\varphi.
\end{equation}

Введемо фундаментальну фундаментальну довжину~--- \emphx{електронний радіус Дебая} $\lambda_{De}$:
\begin{equation}
	\tcbhighmath{\lambda_{De} = \sqrt{\frac{k_{\mathrm B}T_e}{4\pi n_e e^2}}}.
\end{equation}

Тоді рівняння набуває вигляду:
\begin{equation}
	\frac{1}{r^2}\frac{d}{dr}\left(r^2\frac{d\varphi}{dr}\right) - \frac{\varphi}{\lambda_{De}^2} = 0.
\end{equation}

Розв'язком цього рівняння з межевими умовами $\varphi(r) \to \frac{q}{r}$ при $r \to 0$ та $\varphi(r) \to 0$ при $r \to \infty$ є \emphx{потенціал Юкави (або Юкави--Дебая)}:
\begin{equation}
	\varphi(r) = \frac{q}{r} \exp\left(-\frac{r}{\lambda_{De}}\right).
\end{equation}

\textbf{Фізичний зміст результату:} Кулонівський потенціал точкового заряду $1/r$ експоненційно згасає на відстанях порядку $\lambda_{De}$. Тепловий рух прагне розмити хмару зарядів, а електростатичне притягання~--- зібрати її. Баланс цих двох факторів і задає розмір дебаївської «шуби» $\lambda_{De}$.

Якщо іони також мобільні і мають температуру $T_i$, то повна дебаївська довжина $\lambda_D$ враховує обидві компоненти:
\begin{equation}
	\frac{1}{\lambda_D^2} = \frac{1}{\lambda_{De}^2} + \frac{1}{\lambda_{Di}^2} = 4\pi e^2 \left(\frac{n_e}{k_{\mathrm B}T_e} + \frac{Z^2 n_i}{k_{\mathrm B}T_i}\right).
\end{equation}

\begin{figure}[htbp]\centering
    \localinput{potential.tikz}
	\caption{Дебаївське екранування пробного заряду: вакуумний кулонівський потенціал (штрихова лінія) експоненціально послаблюється плазмою (суцільна лінія) на масштабі $\lambda_D$.}\label{fig:debye-screening}
\end{figure}

%% --------------------------------------------------------
\subsection{Дебаївська сфера та критерій колективності}
%% --------------------------------------------------------

\textbf{Навіщо вводиться параметр $N_D$?}
Екранування є статистичним ефектом. Щоб поняття дебаївської «шуби» мало фізичний сенс, усередині сфери радіуса $\lambda_D$ повинно знаходитися багато частинок, які здійснюють екранування.

Об'єм дебаївської сфери дорівнює $V_D = \frac{4\pi}{3}\lambda_D^3$. Число електронів у цій сфері називається \emphx{числом Дебая} $N_D$:
\begin{equation}
	N_D = n_e \cdot \frac{4\pi}{3}\lambda_D^3 = \frac{4\pi}{3} n_e \left(\sqrt{\frac{k_{\mathrm B}T_e}{4\pi n_e e^2}}\right)^3 = \frac{1}{3\sqrt{4\pi}} \frac{(k_{\mathrm B}T_e)^{3/2}}{e^3 n_e^{1/2}}.
\end{equation}

Головна умова існування плазми (\emphx{критерій колективності}):
\begin{equation}
	\tcbhighmath{N_D \gg 1.}
\end{equation}

\begin{SCfigure}[2][h!]\centering
	\localinput{debye_sphere.tikz}
	\caption{Дебаївська сфера радіуса $\lambda_D$. Умова $N_D \gg 1$ гарантує, що екранування створюється колективною дією багатьох частинок.}\label{fig:debye-sphere}
\end{SCfigure}

Порівняємо дебаївську довжину $\lambda_D$ із середньою відстанню між частинками $a \sim n_e^{-1/3}$. Оскільки $N_D \sim (\lambda_D / a)^3$, умова $N_D \gg 1$ еквівалентна вимозі:
\begin{equation}
	\lambda_D \gg a.
\end{equation}

Отже, у фізиці плазми діє \textbf{ієрархія просторових масштабів}:
\begin{equation}
	\tcbhighmath{a \ll \lambda_D \ll L},
\end{equation}
де $a$~--- міжчастинкова відстань, $\lambda_D$~--- масштаб екранування та мікроскопічних збурень, $L$~--- макроскопічний розмір плазмового об'єму (контейнера).

%% --------------------------------------------------------
\subsection{Ідеальна та сильнозв'язана плазма}
%% --------------------------------------------------------

\textbf{Фізична проблема:} Коли плазму можна вважати ідеальним газом?

Введемо безрозмірний \emphx{параметр зв'язку (плазмовий параметр)} $\Gamma$, який описує відношення середньої кулонівської потенціальної енергії взаємодії сусідніх частинок до їхньої теплової кінетичної енергії:
\begin{equation}
	\Gamma = \frac{E_{\text{потенц}}}{E_{\text{кін}}} = \frac{e^2 / a}{k_{\mathrm B}T} = \frac{e^2 n^{1/3}}{k_{\mathrm B}T}.
\end{equation}

Звернемо увагу, що $\Gamma \sim N_D^{-2/3}$. Звідси випливає класифікація:
\begin{enumerate}
	\item \textbf{Слабкозв'язана (ідеальна) плазма ($\Gamma \ll 1 \Leftrightarrow N_D \gg 1$):}
	Тепловий рух домінує над кулонівським притяганням/відштовхуванням сусідніх частинок. Термодинамічно вона описується рівнянням стану ідеального газу $p \approx \sum n_s k_{\mathrm B}T_s$. Більшість космічної та лабораторної плазми є ідеальною.
	\item \textbf{Сильнозв'язана (неідеальна) плазма ($\Gamma \gtrsim 1 \Leftrightarrow N_D \lesssim 1$):}
	Виникає за високих густин та низьких температур (наприклад, у ядрах білих карликів, у плазмі вибуху або рідинних металах). У такій плазмі виникають ближній порядок та кореляційні ефекти, подібні до рідин.
\end{enumerate}

%% --------------------------------------------------------
\subsection{Амбіполярна дифузія}
%% --------------------------------------------------------

\textbf{Фізична проблема:} Електрони мають набагато більшу теплову швидкість, ніж іони ($v_{Te} \gg v_{Ti}$). Якщо виникає градієнт концентрації $\nabla n$, електрони мають дифундувати набагато швидше за іони. Чому ж плазма не розпадається?

Під час розширення плазмової хмари легкі електрони вириваються вперед, залишаючи важкі іони позаду. Це миттєво порушує квазінейтральність і створює внутрішнє роздільне електричне поле $\vect{E}_{\text{амб}}$, напрямлене від іонів до електронів.

Це поле:
\begin{itemize}
	\item Гальмує швидкі електрони;
	\item Прискорює повільні іони.
\end{itemize}

У результаті обидва компоненти вимушені дифундувати з однаковою спільною швидкістю. Цей процес називається \emphx{амбіполярною дифузією}.

Запишемо густини потоків частинок $\vect{\Gamma}_e$ та $\vect{\Gamma}_i$ з урахуванням дифузії та дрейфу в електричному полі:
\begin{align}
	\vect{\Gamma}_e &= -D_e \nabla n_e - \mu_e n_e \vect{E}_{\text{амб}}, \\
	\vect{\Gamma}_i &= -D_i \nabla n_i + \mu_i n_i \vect{E}_{\text{амб}},
\end{align}
де $D_s = \frac{k_{\mathrm B}T_s}{m_s \nu_s}$~--- коефіцієнти дифузії, а $\mu_s = \frac{|q_s|}{m_s \nu_s}$~--- рухливості частинок сорту $s$.

Вимога збереження квазінейтральності за наявності однократно іонізованих частинок ($n_e = n_i = n$) означає рівність потоків: $\vect{\Gamma}_e = \vect{\Gamma}_i = \vect{\Gamma}$. Виражаючи $\vect{E}_{\text{амб}}$ з цієї рівності, отримуємо розрахунок для напрямленого потоку $\vect{\Gamma} = -D_a \nabla n$, де $D_a$~--- \emphx{коефіцієнт амбіполярної дифузії}:
\begin{equation}
	D_a = \frac{D_e \mu_i + D_i \mu_e}{\mu_e + \mu_i}.
\end{equation}

Оскільки $\mu_e \gg \mu_i$ та $D_e / \mu_e = k_{\mathrm B}T_e / e$, то вираз спрощується до:
\begin{equation}
	\tcbhighmath{D_a \approx D_i \left(1 + \frac{T_e}{T_i}\right)}.
\end{equation}

\textbf{Висновок:} Швидкість розширення всієї плазми визначається повільними іонами, але коефіцієнт дифузії збільшується в $(1 + T_e/T_i)$ разів за рахунок виникаючого самоузгодженого поля електронів.

%% ========================================================
\section{Часові масштаби, плазмова частота та хвилі}\label{sec:time-scales}
%% ========================================================

%% --------------------------------------------------------
\subsection{Плазмова (ленгмюрівська) частота}
%% --------------------------------------------------------

\textbf{Методичний місток:} Дебаївський радіус $\lambda_D$ задає \emph{просторову} межу розділення зарядів. Який фундаментальний \emph{часовий} масштаб відповідає реакції плазми на збурення?

Розглянемо спрощену модель однорідного шару плазми. Нехай іони утворюють нерухомий позитивний фон густиною $+e n_0$. Змістимо весь електронний шар як ціле на малу відстань $x$ вздовж осі $OX$ (рис.~\ref{fig:plasma-oscillation}).

\begin{figure}[htbp]\centering
	\begin{subfigure}[t]{0.48\linewidth}\centering
		\localinput{plasma_ions.tikz}
	\end{subfigure}\hfill
	\begin{subfigure}[t]{0.48\linewidth}\centering
		\localinput{condensator.tikz}
	\end{subfigure}
	\caption{Модель плазмових коливань: зміщення електронного шару відносно нерухомих іонів утворює плоский конденсатор із повертаючим електричним полем $\vect{E}$.}\label{fig:plasma-oscillation}
\end{figure}

На межах виникнуть нескомпенсовані поверхневі заряди із густиною $\sigma = e n_0 x$. Це приводить до виникнення електричного поля, аналогічно плоскому конденсатору:
\begin{equation}
	E_x = -4\pi \sigma = -4\pi e n_0 x.
\end{equation}

Рівняння руху для кожного електрона під дією цього повертаючого поля має вигляд:
\begin{equation}
	m_e \frac{d^2 x}{dt^2} = -e E_x = -4\pi n_0 e^2 x \quad \Rightarrow \quad \frac{d^2 x}{dt^2} + \left(\frac{4\pi n_0 e^2}{m_e}\right) x = 0.
\end{equation}

Це рівняння гармонічного осцилятора $\ddot{x} + \omega_{pe}^2 x = 0$, де власна частота коливань дорівнює:
\begin{equation}
	\tcbhighmath{\omega_{pe} = \sqrt{\frac{4\pi n_e e^2}{m_e}}}.
\end{equation}

Ця величина називається \emphx{електронною плазмовою (або ленгмюрівською) частотою}.

\textbf{Фізичний зміст:} $\omega_{pe}$~--- це найшвидший часовий масштаб у плазмі. Вона визначає, як швидко електрони здатні зреагувати та заекранувати будь-яке виникаюче збурення заряду.

Пов'яжемо дебаївську довжину та плазмову частоту через теплову швидкість електронів $v_{Te} = \sqrt{k_{\mathrm B}T_e / m_e}$:
\begin{equation}
	\tcbhighmath{\lambda_{De} = \frac{v_{Te}}{\omega_{pe}}}.
\end{equation}
\textbf{Фізична інтерпретація:} Дебаївський радіус~--- це відстань, яку електрон зі середньою тепловою швидкістю проходить за один період плазмових коливань $\tau \sim \omega_{pe}^{-1}$.

%% --------------------------------------------------------
\subsection{Діелектрична проникність плазми}
%% --------------------------------------------------------

\textbf{Фізична проблема:} Як плазма реагує на зовнішнє високочастотне електромагнітне випромінювання (наприклад, радіохвилю або лазерний промінь)?

Розглянемо плоску монохроматичну хвилю $E(t) = E_0 e^{-i\omega t}$, що поширюється в плазмі. Оскільки маса іонів велика, розглянемо лише рух електронів під дією електричного поля хвилі (нехтуючи магнітною силою Лоренца $v/c \ll 1$ та зіткненнями):
\begin{equation}
	m_e \frac{d^2 \vect{r}}{dt^2} = -e \vect{E}_0 e^{-i\omega t} \quad \Rightarrow \quad \vect{r}(t) = \frac{e}{m_e \omega^2}\vect{E}(t).
\end{equation}

Зміщення електронів створює макроскопічний вектор поляризації середовища $\vect{P}$:
\begin{equation}
	\vect{P} = -e n_e \vect{r}(t) = -\frac{n_e e^2}{m_e \omega^2}\vect{E}.
\end{equation}

За визначенням електродинаміки суцільних середовищ, вектор електричної індукції $\vect{D} = \vect{E} + 4\pi\vect{P} = \varepsilon(\omega)\vect{E}$. Звідси знаходимо \emphx{діелектричну проникність холодної беззіткнювальної плазми}:
\begin{equation}
	\varepsilon(\omega) = 1 - \frac{4\pi n_e e^2}{m_e \omega^2} \quad \Rightarrow \quad \tcbhighmath{\varepsilon(\omega) = 1 - \frac{\omega_{pe}^2}{\omega^2}}.
\end{equation}

\textbf{Фізичний аналіз формули:}
\begin{itemize}
	\item На відміну від звичайних діелектриків, де $\varepsilon > 1$, у плазмі $\varepsilon(\omega) < 1$.
	\item При високочастотному полі ($\omega \gg \omega_{pe}$) електрони за інерцією не встигають зміщуватися, $\varepsilon \to 1$, і плазма поводиться як вакуум.
	\item При $\omega < \omega_{pe}$ діелектрична проникність стає \textbf{від'ємною} ($\varepsilon < 0$).
\end{itemize}

%% --------------------------------------------------------
\subsection{Закон дисперсії та відбиття хвиль від плазми}
%% --------------------------------------------------------

\textbf{Фізична проблема:} Як від'ємне значення $\varepsilon(\omega)$ впливає на поширення електромагнітних хвиль?

Запишемо хвильове рівняння, що випливає з рівнянь Максвелла для монохроматичного поля $\vect{E} \propto e^{i(\vect{k}\cdot\vect{r} - \omega t)}$:
\begin{equation}
	k^2 = \frac{\omega^2}{c^2}\varepsilon(\omega) = \frac{\omega^2}{c^2}\left(1 - \frac{\omega_{pe}^2}{\omega^2}\right).
\end{equation}

Перегрупувавши члени, отримуємо \emphx{закон дисперсії поперечних хвиль у плазмі}:
\begin{equation}
	\tcbhighmath{\omega^2 = \omega_{pe}^2 + c^2 k^2}.
\end{equation}

Проаналізуємо цей закон (рис.~\ref{fig:plasma-dispersion-graph}):

\begin{figure}[htbp]\centering
    \localinput{dispersion.tikz}
	\caption{Дисперсійна крива електромагнітної хвилі в плазмі. При $\omega < \omega_{pe}$ хвильовий вектор $k$ є уявним~--- хвиля проникає лише на глибину скінченного скін-шару і повністю відбивається.}\label{fig:plasma-dispersion-graph}
\end{figure}

\begin{itemize}
	\item \textbf{Режим прозорості ($\omega > \omega_{pe}$):}
	Хвильове число $k = \frac{\sqrt{\omega^2 - \omega_{pe}^2}}{c}$ є дійсним. Хвиля вільно поширюється у плазмі зі фазовою швидкістю $v_{\text{ф}} = \frac{\omega}{k} = \frac{c}{\sqrt{1 - \omega_{pe}^2/\omega^2}} > c$ та груповою швидкістю $v_{\text{гр}} = \frac{d\omega}{dk} = c\sqrt{1 - \frac{\omega_{pe}^2}{\omega^2}} < c$.

	\item \textbf{Режим відбиття / Непрозорості ($\omega < \omega_{pe}$):}
	Величина $k^2 < 0$, тобто хвильове число є чисто уявним: $k = i\kappa = i\frac{\sqrt{\omega_{pe}^2 - \omega^2}}{c}$.

	Підставляючи це у вираз для амплітуди хвилі $E \propto e^{i k x} = e^{-\kappa x}$, бачимо, що амплітуда хвилі \textbf{експоненційно згасає} вглиб плазми. Хвиля не може поширюватися і повністю відбивається від межі плазми (коефіцієнт відбиття за інтенсивністю $R = 1$). Глибина проникнення (скін-шар) дорівнює:
	\begin{equation}
		\delta = \frac{1}{\kappa} = \frac{c}{\sqrt{\omega_{pe}^2 - \omega^2}} \xrightarrow{\omega \ll \omega_{pe}} \frac{c}{\omega_{pe}} \equiv d_e \quad (\text{безінерційний скін-шар}).
	\end{equation}
\end{itemize}

\textbf{Практичне застосування (Іоносферний радіозв'язок):}
Верхній шар атмосфери Землі~--- іоносфера~--- являє собою плазму з концентрацією електронів $n_e \sim 10^5$--$10^6~\mathrm{см}^{-3}$. Відповідна плазмова частота становить:
\begin{equation}
	\nu_{pe} = \frac{\omega_{pe}}{2\pi} \approx 9\sqrt{n_e\,[\mathrm{см}^{-3}]}~\mathrm{kHz} \sim 3\text{--}9~\mathrm{MHz}.
\end{equation}
Радіохвилі короткохвильового діапазону (КХ) з частотою $\nu < \nu_{pe}$ відбиваються від іоносфери як від дзеркала, що дає змогу здійснювати закритичний заризонтний радіозв'язок навколо всієї земної кулі. Супутникові ж сигнали (GPS, Wi-Fi, ТБ) мають частоти у діапазоні $\mathrm{GHz} \gg \nu_{pe}$, тому вільно проходять крізь іоносферу на Землю.

%% ========================================================
\section{Транспортні явища та провідність Спітцера}\label{sec:transport}
%% ========================================================

%% --------------------------------------------------------
\subsection{Кулонівські зіткнення та кулонівський логарифм}
%% --------------------------------------------------------

\textbf{Методична проблема:} Чому парні зіткнення в плазмі принципово відрізняються від зіткнень більярдних куль?

Коли електрон пролітає повз іон, він відчуває дію кулонівської сили $F \propto 1/r^2$. Відхилення на великий кут ($\theta \approx 90^\circ$) відбувається лише за умови дуже близького прольоту з прицільним параметром $b_0$:
\begin{equation}
	\frac{Z e^2}{b_0} \approx m_e v^2 \quad \Rightarrow \quad b_0 = \frac{Z e^2}{m_e v^2}.
\end{equation}

Однак кількість «далеких» прольотів із прицільним параметром $b > b_0$ є значно більшою. При кожному далекому прольоті електрон відхиляється на мізерний кут, але сумарна дія тисяч таких розсіянь накопичується дифузійним чином.

Виявляється, що \textbf{сумарний ефект накопичення малих відхилень у $8\ln\Lambda$ разів ефективніший, ніж поодинокі лобові зіткнення!}

Параметр $\ln\Lambda$ називається \emphx{кулонівським логарифмом}:
\begin{equation}
	\Lambda = \frac{b_{\max}}{b_{\min}}.
\end{equation}
Фізичні межі інтегрування дорівнюють:
\begin{itemize}
	\item Максимальний прицільний параметр $b_{\max} = \lambda_D$ (на більших відстанях кулонівське поле повністю закрите дебаївським екрануванням);
	\item Мінімальний прицільний параметр $b_{\min} = \max(b_0, \lambda_{\text{де Бройля}})$, де $b_0 = \frac{Z e^2}{3 k_{\mathrm B}T_e}$.
\end{itemize}

У термоядерній плазмі ratio $b_{\max}/b_{\min} \sim 10^6$, тому кулонівський логарифм слабко залежить від конкретних параметрів і зазвичай лежить у межах:
\begin{equation}
	\tcbhighmath{\ln\Lambda \approx 10\text{--}20}.
\end{equation}

%% --------------------------------------------------------
\subsection{Провідність Спітцера}
%% --------------------------------------------------------

\textbf{Фізична проблема:} Як залежить електрична провідність повністю іонізованої плазми від її температури?

Запишемо дрейфову швидкість електронів $v_d$ у постійному електричному полі $E$ за наявності сили кулонівського тертя $m_e v_d / \tau_{ei}$:
\begin{equation}
	-e E - \frac{m_e v_d}{\tau_{ei}} = 0 \quad \Rightarrow \quad v_d = -\frac{e \tau_{ei}}{m_e} E,
\end{equation}
де $\tau_{ei}$~--- ефективний час між кулонівськими зіткненнями електронів з іонами.

Густина струму $j = -e n_e v_d = \sigma E$, звідки класична провідність становить:
\begin{equation}
	\sigma = \frac{n_e e^2 \tau_{ei}}{m_e}.
\end{equation}

Підставимо у цю формулу час вільного пробігу з урахуванням кулонівського логарифма $\tau_{ei} = \frac{1}{n_i v_{Te} \sigma_{\text{кулон}}} \sim \frac{(k_{\mathrm B}T_e)^{3/2} \sqrt{m_e}}{Z n_i e^4 \ln\Lambda}$:
\begin{equation}
	\tcbhighmath{\sigma_{\text{Spitzer}} = \gamma_E \frac{(k_{\mathrm B}T_e)^{3/2}}{Z e^2 m_e^{1/2} \ln\Lambda} \propto \frac{T_e^{3/2}}{Z \ln\Lambda}},
\end{equation}
де $\gamma_E \approx 0.58$ (для $Z=1$)~--- поправковий коефіцієнт Лаймана Спітцера, що враховує електрон-електронні зіткнення.

\textbf{Дивовижний фізичний висновок:}
\begin{itemize}
	\item \textbf{Незалежність від густини:} Провідність повністю іонізованої плазми \emph{не залежить від концентрації частинок} $n_e$. При збільшенні $n_e$ зростає число носіїв заряду, але пропорційно зростає й частота зіткнень, що повністю компенсує ефект.
	\item \textbf{Залежність від температури ($T_e^{3/2}$):} Чим гарячішою стає плазма, тим швидше рухаються електрони, тим важче іонам їх відхилити, і тим \textbf{меншим} стає ефективний опір!
	За високих температур ($T_e > 1~\mathrm{keV}$) провідність плазми перевищує провідність надчистої міді при кімнатній температурі. Це унеможливлює нагрівання високотемпературної плазми звичайним омічним струмом (омічний нагрів згасає при $T_e \sim 2\text{--}3~\mathrm{keV}$).
\end{itemize}

%% ========================================================
\section{Магнітне утримання та магнітна гідродинаміка}\label{sec:magnetic-confinement}
%% ========================================================

%% --------------------------------------------------------
\subsection{Рух зарядженої частинки в магнітному полі}
%% --------------------------------------------------------

\textbf{Методичний місток:} Як закрутити нейтральний газ, ми знаємо (стінки камери). Як утримати гарячу плазму, яка розплавить будь-яку матеріальну стінку? За допомогою магнітного поля!

На заряджену частинку $q$, що рухається зі швидкістю $\vect{v}$ у магнітному полі $\vect{B}$, діє сила Лоренца:
\begin{equation}
	\vect{F} = \frac{q}{c}\vect{v}\times\vect{B}.
\end{equation}

Розкладемо швидкість на поздовжню $v_\parallel$ та перпендикулярну $v_\perp$ відносно напрямку $\vect{B}$. Поздовжній рух є вільним, а в перпендикулярній площині сила Лоренца виступає як доцентрова:
\begin{equation}
	\frac{m v_\perp^2}{r_L} = \frac{|q|}{c} v_\perp B.
\end{equation}

Звідси визначаються два ключові параметри обертання:
\begin{align}
	&\text{\emphx{Циклотронна (ларморівська) частота:}} \quad \tcbhighmath{\omega_c = \frac{|q|B}{m c}}, \label{eq:cyclotron-freq} \\
	&\text{\emphx{Ларморівський радіус (радіус обертання):}} \quad \tcbhighmath{r_L = \frac{v_\perp}{\omega_c} = \frac{m c v_\perp}{|q|B}}. \label{eq:larmor-radius}
\end{align}

Частинка намотується на силову лінію магнітного поля, рухаючись по спіралі (рис.~\ref{fig:larmor-motion}).

\begin{figure}[htbp]\centering
	\begin{tikzpicture}[>=latex,scale=0.85]
		\draw[->] (-4.2,0)--(4.4,0) node[right] {$z\parallel\vect{B}$};
		\draw[->] (0,-2.2)--(0,2.5) node[above] {$y$};
		\draw[thick,domain=-3.7:3.7,samples=180] plot (\x,{1.15*sin(220*\x/360)});
		\node at (2.3,1.55) {$r_L$};
		\node at (0,-1.75) {$\displaystyle\omega_c=\frac{|q|B}{mc}$};
	\end{tikzpicture}
	\caption{Гвинтовий рух частинки вздовж магнітного поля. Магнітне поле позбавляє частинку свободи руху в двох поперечних напрямках.}\label{fig:larmor-motion}
\end{figure}

Оскільки $m_e \ll m_i$, електронний ларморівський радіус набагато менший за іонний ($r_{Le} \ll r_{Li}$). Магнітне поле обмежує поперечний рух частинок розміром $r_L \ll L$, що робить можливим магнітне утримання.

%% --------------------------------------------------------
\subsection{Магнітний момент та принцип магнітного дзеркала}
%% --------------------------------------------------------

Обертання частинки по колу радіуса $r_L$ є замкненим мікроскопічним струмом $I = \frac{|q|\omega_c}{2\pi c}$. Цей струм створює \emphx{магнітний момент} $\mu$:
\begin{equation}
	\mu = \frac{I \cdot S}{c} = \frac{|q|\omega_c}{2\pi c} \cdot (\pi r_L^2) = \frac{m v_\perp^2}{2B} = \frac{E_{\perp}}{B}.
\end{equation}

При повільній зміні магнітного поля в просторі або часі (порівняно з періодом $\omega_c^{-1}$) величина $\mu$ зберігається і є \emphx{першим адіабатичним інваріантом}:
\begin{equation}
	\tcbhighmath{\mu = \frac{m v_\perp^2}{2B} \approx \text{const}.}
\end{equation}

Розглянемо рух частинки в \textbf{неоднорідному магнітному полі}, яке посилюється вздовж силової лінії (наприклад, на крайках пастки, де $B$ зростає від $B_{\min}$ до $B_{\max}$).

Оскільки магнітне поле не виконує роботи над зарядженою частинкою, повна кінетична енергія частинки зберігається:
\begin{equation}
	E = E_\parallel + E_\perp = \frac{m v_\parallel^2}{2} + \mu B = \text{const}.
\end{equation}

Під час руху частинки в область сильнішого поля ($B \uparrow$):
\begin{itemize}
	\item Щоб зберегти $\mu = E_\perp / B$, поперечна енергія мусить зростати ($E_\perp \uparrow$);
	\item Щоб зберегти повну енергію $E$, поздовжня енергія повинна зменшуватися ($E_\parallel \downarrow$).
\end{itemize}

У точці, де $E_\parallel$ стає рівною нулю ($v_\parallel = 0$), частинка зупиняється і відштовхується назад в область слабкого поля. Цей ефект називається \emphx{магнітним дзеркалом (пробкою)} (рис.~\ref{fig:magnetic-mirror}).

\begin{figure}[htbp]\centering
    \localinput{mirror.tikz}
	\caption{Принцип магнітного дзеркала: зростання $B$ викликає збільшення $v_\perp$ за рахунок $v_\parallel$, викликаючи відбиття частинки від області сильного поля.}\label{fig:magnetic-mirror}
\end{figure}

%% --------------------------------------------------------
\subsection{Конус втрат у відкритих пастках}
%% --------------------------------------------------------

\textbf{Фізична проблема:} Чи всі частинки відбиваються магнітним дзеркалом?

Ні. Нехай у центрі пастки (де поле $B_{\min}$) частинка має пітч-кут $\theta$ між вектором швидкості та силовою лінією ($\sin\theta = v_\perp / v$). Повна умова відбиття частинки до досягнення максимуму поля $B_{\max}$ має вигляд:
\begin{equation}
	\tcbhighmath{\sin^2\theta \ge \frac{B_{\min}}{B_{\max}} = \frac{1}{R_m}},
\end{equation}
де $R_m = B_{\max}/B_{\min} > 1$~--- дзеркальне відношення.

Якщо кут частинки занадто малий ($\theta < \theta_{\text{lc}}$), вона пролітає крізь дзеркало і залишає пастку. Ці частинки утворюють \emphx{конус втрат} у просторі швидкостей з кутом розкриття:
\begin{equation}
	\theta_{\text{lc}} = \arcsin\left(\frac{1}{\sqrt{R_m}}\right).
\end{equation}
Втрати частинок через конус втрат є головним недоліком відкритих магнітних пасток.

%% --------------------------------------------------------
\subsection{Параметр \texorpdfstring{$\beta$}{beta} та МГД-рівновага}
%% --------------------------------------------------------

\textbf{Методичний місток:} Як оцінити ефективність утримання плазми магнітним полем?

Запишемо газодинамічний тиск плазми $p = \sum n_s k_{\mathrm B}T_s = n_e k_{\mathrm B}T_e + n_i k_{\mathrm B}T_i$.
З іншого боку, густина енергії магнітного поля створює ефективний \emphx{магнітний тиск} $p_B = \frac{B^2}{8\pi}$.

Відношення цих двох величин називається \emphx{параметром бета ($\beta$)}:
\begin{equation}
	\tcbhighmath{\beta = \frac{p}{B^2 / (8\pi)} = \frac{8\pi p}{B^2}}.
\end{equation}

Фізичний зміст $\beta$:
\begin{itemize}
	\item $\beta \ll 1$~--- магнітне поле повністю домінує, плазма є «пасивною» і не деформує сильні зовнішні магнітні лінії (режим сучасних токамаків, де $\beta \sim 1$--$5\%$).
	\item $\beta \sim 1$~--- тиск плазми порівнянний із магнітним тиском; плазма здатна сильно витісняти магнітне поле (високоефективний режим, але схильний до нестійкостей).
\end{itemize}

Запишемо **основне рівняння МГД-рівноваги** ($\vect{v} = 0$) з урахуванням амперової сили $\frac{1}{c}\vect{j}\times\vect{B}$:
\begin{equation}
	\nabla p = \frac{1}{c}\vect{j}\times\vect{B}.
\end{equation}

Застосувавши рівняння Максвелла $\rot\vect{B} = \frac{4\pi}{c}\vect{j}$, перетворимо амперову силу:
\begin{equation}
	\frac{1}{c}\vect{j}\times\vect{B} = \frac{1}{4\pi}(\rot\vect{B})\times\vect{B} = \frac{1}{4\pi}(\vect{B}\cdot\nabla)\vect{B} - \nabla\left(\frac{B^2}{8\pi}\right).
\end{equation}

Тоді рівняння рівноваги набуває яскравого фізичного вигляду:
\begin{equation}
	\tcbhighmath{\nabla \left(p + \frac{B^2}{8\pi}\right) = \frac{(\vect{B}\cdot\nabla)\vect{B}}{4\pi}}.
\end{equation}

\textbf{Фізична інтерпретація:}
Сума газокінетичного тиску $p$ та магнітного тиску $B^2/8\pi$ зрівноважується силою натягу магнітних силових ліній $\frac{(\vect{B}\cdot\nabla)\vect{B}}{4\pi}$ (сила, що прагне розпрямити викривлені силові лінії).

%% --------------------------------------------------------
\subsection{Пінч-ефект, умова Беннета та нестійкості}
%% --------------------------------------------------------

\textbf{Фізична ідея:} Струм, що тече через плазмовий циліндр, створює власне охоплююче азимутальне магнітне поле $B_\theta$. Взаємодія струму з власним полем створює амперову силу, напрямлену до осі, яка стискає плазмовий шнур. Це явище називається \emphx{$z$-пінчем}.

\begin{figure}[h!]\centering
	% \localinput{z-pinch-equilibrium.tikz}
	\caption{Схема $z$-пінча: струм $I$, що тече вздовж осі $Z$, створює азимутальне магнітне поле $B_\theta$, яке стискає шнур газового розряду.}\label{pic:z-pinch-equilibrium}
\end{figure}

Знайдемо умову рівноваги для циліндричного пінча радіуса $R$. Інтегруючи рівняння МГД-рівноваги по перетину шнура, отримуємо класичну \emphx{умову рівноваги Беннета}:
\begin{equation}
	\tcbhighmath{2 N_1 k_{\mathrm B}(T_e + T_i) = \frac{I^2}{c^2}},
\end{equation}
де $N_1 = \int_0^R 2\pi r n(r) dr$~--- погонна кількість частинок на одиницю довжини шнура, $I$~--- повний струм у СГС.

\textbf{Фізичний зміст умови Беннета:} Щоб утримати плазмовий шнур із заданим числом частинок $N_1$ та температурою $T$, необхідно пропустити струм $I \propto \sqrt{N_1 T}$.

\textbf{Нестійкості пінча:}
На жаль, простий $z$-пінч є МГД-нестійким. Виділяють дві основні моди нестійкостей:
\begin{itemize}
	\item \textbf{Перетяжка ($m=0$):} Випадкове локальне звуження радіуса шнура $R \to R - \delta R$ приводить до локального зростання поля $B_\theta(R) \propto 1/R$ та магнітного тиску $\propto 1/R^2$. Поле перетискає шнур усе сильніше, аж до його повного розриву.
	\item \textbf{Згин ($m=1$):} Випадкове викривлення шнура приводить до густішого розташування силових ліній з увігнутого боку. Магнітний тиск там зростає і викидає вигин назовні.
\end{itemize}

Для стабілізації пінча всередину шнура «вморожують» сильне поздовжнє магнітне поле $B_z$, яке чинить опір стисканню та вигинам.

%% ========================================================
\section{Керований термоядерний синтез (КТС)}\label{sec:fusion}
%% ========================================================

%% --------------------------------------------------------
\subsection{Фізичні основи та реакція D-T}
%% --------------------------------------------------------

Термоядерний синтез~--- це процес об'єднання легких атомних ядер у важчі, який супроводжується виділенням величезної кількості енергії завдяки дефекту маси. Найбільш досяжною для лабораторної реалізації є реакція між ізотопами водню~--- дейтерієм ($\mathrm{D}$) та тритієм ($\mathrm{T}$):
\begin{equation}
	\tcbhighmath{{}_1^2\mathrm{H} + {}_1^3\mathrm{H} \longrightarrow {}_2^4\mathrm{He}\,(3.5~\mathrm{MeV}) + {}_0^1\mathrm{n}\,(14.1~\mathrm{MeV})}.
\end{equation}

Загальне енерговиділення становить $17.6~\mathrm{MeV}$ на один акт реакції. Альфа-частинки (${}_2^4\mathrm{He}$) затримаються у плазмі магнітним полем, підтримуючи її високу температуру, а високоенергетичні нейтрони залишають плазму і поглинаються в бланкеті реактора, виділяючи тепло для виробництва електроенергії.

%% --------------------------------------------------------
\subsection{Критерій Лоусона}
%% --------------------------------------------------------

\textbf{Фізична проблема:} За яких параметрів плазми реактор вироблятиме більше енергії, ніж витрачається на його нагрівання та компенсацію втрат?

Енергетичний баланс вимагає, щоб виділена термоядерна потужність перевищувала втрати на теплопровідність та випромінювання. Це дає \emphx{критерій Лоусона}:
\begin{equation}
	\tcbhighmath{n_e \cdot \tau_E \ge 10^{14}~\frac{\mathrm{s}}{\mathrm{см}^3} \quad (\text{для D-T реакції при } T \sim 10\text{--}20~\mathrm{keV})},
\end{equation}
де $n_e$~--- концентрація плазми, $\tau_E$~--- \emphx{енергетичний час утримання} (час, за який плазма втрачає свою теплову енергію за рахунок теплопровідності).

Існує два основні шляхи досягнення критерію Лоусона:
\begin{itemize}
	\item \textbf{Магнітне утримання:} Поміркована густина ($n_e \sim 10^{14}~\mathrm{см}^{-3}$), але тривалий час утримання ($\tau_E \sim 1\text{--}10~\mathrm{с}$).
	\item \textbf{Інерційне утримання:} Величезна густина при надвисокому стисненні ($n_e \sim 10^{25}~\mathrm{см}^{-3}$), але надкороткий час утримання за рахунок власної інерції мікромішені ($\tau_E \sim 10^{-10}~\mathrm{с}$).
\end{itemize}

%% --------------------------------------------------------
\subsection{Тороїдальні магнітні пастки: Токамаки vs Стеларатори}
%% --------------------------------------------------------

Щоб уникнути втрат частинок через відкриті торці, магнітні силові лінії замкають у тороїдальну конфігурацію. Проте простий тороїдальний загин приводить до неоднорідності поля $B \propto 1/R$, що викликає градієнтний дрейф електронів та іонів у протилежних напрямках та швидкий викид плазми на стінку.

Щоб компенсувати дрейф, необхідна \textbf{гвинтова обкрутка} силових ліній навколо тора. За способом створення гвинтового поля замкнені тороїдальні системи поділяють на два класи:

\begin{enumerate}
	\item \textbf{Токамак (Тороїдальна камера з магнітними котушками):}
	Гвинтове поле створюється комбінацією сильного тороїдального поля від зовнішніх котушок та попольного поля, яке створюється \textbf{великим поздовжнім струмом}, що індукується безпосередньо в плазмовому шнурі за допомогою трансформаторного сердечника.

	\item \textbf{Стеларатор:}
	Уся гвинтова конфігурація магнітного поля створюється повністю \textbf{зовнішніми складнопрофільними 3D-котушками}. Вона не вимагає протікання великого струму через плазму, що забезпечує стаціонарний (безімпульсний) режим роботи.
\end{enumerate}

%% --------------------------------------------------------
\subsection{Сучасний стан термоядерних досліджень}
%% --------------------------------------------------------

Галузь керованого термоядерного синтезу наразі переживає бурхливий перехід від фундаментальної фізики до інженерної реалізації:

\begin{itemize}
	\item \textbf{NIF (National Ignition Facility, США):} У грудні 2022 року на лазерній установці інерційного синтезу NIF уперше в історії людства було досягнуто \emphx{термоядерного запалювання (Ignition)}: лазерна енергія $2.05~\mathrm{MJ}$ поглинулася мішенню, виділивши $3.15~\mathrm{MJ}$ термоядерної енергії ($Q \approx 1.5$).

	\item \textbf{ITER (Міжнародний термоядерний експериментальний реактор, Кадараш, Франція):} Найбільший у світі токамак із надпровідними магнітами, що будується зусиллями міжнародної консорціум-коаліції. Мета ITER~--- продемонструвати $Q \ge 10$ (отримання $500~\mathrm{MW}$ термоядерної потужності при $50~\mathrm{MW}$ вкладеної).

	\item \textbf{Wendelstein 7-X (Грайфсвальд, Німеччина):} Найбільший у світі оптимізований надпровідний стеларатор, який успішно демонструє стабільне утримання плазми високої густини в стаціонарному режимі без розривних нестійкостей.

	\item \textbf{SPARC (MIT / Commonwealth Fusion Systems, США):} Проєкт компактного токамака нового покоління, який використовує високотемпературні надпровідники другого покоління (HTS на основі REBCO). Вони дозволяють створити рекордне магнітне поле $B \approx 12~\mathrm{T}$, що суттєво зменшує розміри реактора і прискорює його комерціалізацію.
\end{itemize}

%% ========================================================
\section*{Підсумки}
%% ========================================================

\begin{itemize}
	\item \textbf{Плазма}~--- це квазінейтральний іонізований газ, поведінка якого визначається колективними кулонівськими ефектами.
	\item \textbf{Фундаментальна ієрархія просторових масштабів:}
	\begin{equation*}
		a \ll \lambda_D \ll L,
	\end{equation*}
	де $\lambda_D = \sqrt{\frac{k_{\mathrm B}T_e}{4\pi n_e e^2}}$~--- радіус дебаївського екранування.
	\item \textbf{Умова колективності:} $N_D = \frac{4\pi}{3}n_e\lambda_D^3 \gg 1$.
	\item \textbf{Часовий масштаб:} Власна частота плазмових коливань $\omega_{pe} = \sqrt{\frac{4\pi n_e e^2}{m_e}}$ визначає динамічну відповідь електронного компонента.
	\item \textbf{Прозорість плазми:} Діелектрична проникність $\varepsilon(\omega) = 1 - \omega_{pe}^2/\omega^2$. Електромагнітні хвилі з частотою $\omega < \omega_{pe}$ відбиваються від плазми.
	\item \textbf{Провідність Спітцера:} $\sigma \propto T_e^{3/2} / (Z \ln\Lambda)$ зростає з температурою і не залежить від густини.
	\item \textbf{Магнітне утримання:} Частинки обертаються з циклотронною частотою $\omega_c = \frac{|q|B}{mc}$ та ларморівським радіусом $r_L = \frac{mc v_\perp}{|q|B}$. Магнітний момент $\mu = \frac{mv_\perp^2}{2B}$ є адіабатичним інваріантом.
	\item \textbf{Параметр $\beta$:} $\beta = \frac{8\pi p}{B^2}$ характеризує співвідношення між газовим та магнітним тисками.
	\item \textbf{Критерій Лоусона для КТС:} $n \tau_E \ge 10^{14}~\mathrm{s}/\mathrm{см}^3$ (для D-T реакції).
\end{itemize}

%% ========================================================
\section*{Задачі}
%% ========================================================

\begin{problem}\label{prb:PlasmaDebye}
	\emphx{Дебаївське екранування.} Лабораторна плазма має електронну концентрацію $n_e = 10^{12}~\mathrm{см}^{-3}$ і температуру електронів $T_e = 2~\mathrm{eV}$. Знайдіть електронну дебаївську довжину $\lambda_{De}$ та число електронів у дебаївській сфері $N_D$.
\end{problem}

\begin{problem}\label{prb:PlasmaCollective}
	\emphx{Умова колективності.} Для розрідженої плазми сонячної корони з концентрацією $n_e = 10^{8}~\mathrm{см}^{-3}$ і температурою $T_e = 100~\mathrm{eV}$ оцініть середню відстань між частинками $a \sim n_e^{-1/3}$ та дебаївську довжину $\lambda_D$. Перевірте виконання умови колективності.
\end{problem}

\begin{problem}\label{prb:PlasmaFrequency}
	\emphx{Плазмова частота та радіозв'язок.} Концентрація електронів в іоносферній хмарі дорівнює $n_e = 2.5 \times 10^{6}~\mathrm{см}^{-3}$. Знайдіть плазмову частоту $\nu_{pe} = \omega_{pe}/(2\pi)$. Чи зможе крізь таку хмару пройти радіосигнал із частотою $\nu = 10~\mathrm{MHz}$?
\end{problem}

\begin{problem}\label{prb:PlasmaConductivity}
	\emphx{Провідність Спітцера.} Термоядерна плазма в токамаку нагрівається від $T_{e1} = 100~\mathrm{eV}$ до $T_{e2} = 4~\mathrm{keV}$. У скільки разів зросте її питома електрична провідність $\sigma$?
\end{problem}

\begin{problem}\label{prb:PlasmaLarmor}
	\emphx{Ларморівський радіус.} Електрон та протон мають однакову кінетичну енергію $E = 1~\mathrm{keV}$ у магнітному полі $B = 20~\mathrm{kG}$ ($2~\mathrm{T}$). Знайдіть і порівняйте їхні ларморівські радіуси, якщо швидкості частинок перпендикулярні до вектора поля.
\end{problem}

\begin{problem}\label{prb:PlasmaMirror}
	\emphx{Конус втрат.} У пробкотроні (відкритій магнітній пастці) магнітне поле в центрі дорівнює $B_{\min} = 10~\mathrm{kG}$, а на кінцях~--- $B_{\max} = 40~\mathrm{kG}$. Обчисліть критичний кут конуса втрат $\theta_{\text{lc}}$. Які частинки вилетять із пастки?
\end{problem}

\begin{problem}\label{prb:BennettPinch}
	\emphx{Умова Беннета.} Знайдіть струм $I$, необхідний для утримання в $z$-пінчі дейтерієвої плазми з погонною густиною $N_1 = 10^{17}~\mathrm{см}^{-1}$ при температурі $T_e = T_i = 1~\mathrm{keV}$.
\end{problem}